% Chapitre 3 : Outils et exemples pratiques

\section{Outils et exemples pratiques}

\subsection{Installation et configuration de BRAT}

\subsubsection{Prérequis système}
\begin{itemize}
    \item Serveur web (Apache, Nginx)
    \item Python 2.7 ou 3.x
    \item Navigateur web moderne (Chrome, Firefox, Safari)
\end{itemize}

\subsubsection{Installation}
\begin{enumerate}
    \item Télécharger BRAT depuis le site officiel
    \item Décompresser l'archive dans le répertoire web
    \item Configurer les permissions d'accès
    \item Tester l'installation avec les données d'exemple
\end{enumerate}

\subsection{Configuration d'un projet}

\subsubsection{Structure des dossiers}
\texttt{project/}
\begin{itemize}
    \item \texttt{annotation.conf} : Configuration du schéma d'annotation
    \item \texttt{visual.conf} : Configuration de l'affichage
    \item \texttt{tools.conf} : Outils d'annotation automatique
    \item \texttt{kb\_shortcuts.conf} : Raccourcis clavier personnalisés
\end{itemize}

\subsubsection{Exemple de configuration}
\begin{verbatim}
[entities]
Person
Organization  
Location
Date

[relations]
Lives_in    Arg1:Person, Arg2:Location
Works_for   Arg1:Person, Arg2:Organization
\end{verbatim}

\subsection{Tutoriel d'annotation pas à pas}

\subsubsection{Annotation d'entités}
\begin{enumerate}
    \item Sélectionner le texte à annoter
    \item Choisir le type d'entité dans le menu contextuel
    \item Valider l'annotation
    \item Vérifier la cohérence avec les annotations existantes
\end{enumerate}

\subsubsection{Annotation de relations}
\begin{enumerate}
    \item Sélectionner la première entité (source)
    \item Glisser vers la seconde entité (cible)
    \item Choisir le type de relation
    \item Confirmer la création de la relation
\end{enumerate}

\subsection{Outils d'analyse et de visualisation}

\subsubsection{Statistiques d'annotation}
\begin{itemize}
    \item Nombre d'entités par type
    \item Distribution des relations
    \item Métriques de progression
\end{itemize}

\subsubsection{Export des données}
\begin{itemize}
    \item Format standoff (natif BRAT)
    \item Export JSON pour intégration
    \item Conversion vers d'autres formats (CoNLL, XML)
\end{itemize}

\subsection{Intégration avec d'autres outils}

\subsubsection{Pipeline de traitement}
\begin{itemize}
    \item Pré-traitement automatique
    \item Annotation semi-automatique
    \item Post-traitement et validation
\end{itemize}

\subsubsection{APIs et scripts}
\begin{itemize}
    \item Scripts Python pour l'import/export
    \item API REST pour l'intégration
    \item Outils de conversion de formats
\end{itemize}

\newpage